\origin{ai}

\subsection{Modeled tAI survival matrix for the $B^{*}_{Q6}$ residual coordinates}
\label{sec:bioreality-bstar-q6-modeled-tai-survival}

\paragraph{Finite surface.}
The record considered here is a finite matrix whose rows are the nine $B^{*}_{Q6}$ synonymous residual coordinates and whose columns are bounded translation readouts derived from organism-scoped codon adaptation weights.  The row set is
$$
\begin{aligned}
R_{Q6}=\{&\mathrm{Arg\_AGR},\mathrm{Ile\_AUA},\mathrm{K\_AAA},\mathrm{Leu\_CUN\_vs\_UUR},\mathrm{Leu\_UUA\_vs\_UUG},\\
&\mathrm{Ser\_UCA\_vs\_UCG},\mathrm{Ser\_UCR\_vs\_AGY},\mathrm{Thr\_ACR\_vs\_ACY},f_{3}\text{-}\mathrm{stress}\}.
\end{aligned}
$$
The codon ambient space contains $61$ sense codons.  The readout used in this finite run is modeled $tAI$ codon weight $w_{c}$ computed from GtRNAdb tRNA gene-copy counts together with the dos Reis wobble $s$-values.  This is a translation-layer annotation surface, not an observed ribosome footprint surface.  The run therefore tests bounded modeled coupling between residual codon coordinates and adaptation-weight readbacks; it does not by itself test translation execution inside the ribosome.

\paragraph{Organism rows and readout completeness.}
The finite organism panel contains $3$ organisms: \emph{Escherichia coli} K-12, \emph{Homo sapiens}, and \emph{Saccharomyces cerevisiae}.  The local GtRNAdb payload covers all $20$ standard amino-acid families in each organism.  For \emph{E. coli} K-12 the FASTA payload has $89$ headers, of which $87$ match the standard tRNA-family parsing and $2$ remain unmatched.  For human the payload has $432$ headers, all $432$ matched.  For yeast the payload has $275$ headers, all $275$ matched.  The modeled $tAI$ construction returns all $61$ sense-codon weights in all three organisms, with zero raw-$W$ fallback count equal to $0$ for \emph{E. coli} K-12, human, and yeast.  Thus the finite row construction is complete at the codon-weight level: missing tRNA-family coverage is not the reason for refusing a stronger biological interpretation.

\paragraph{Matrix statistic.}
The bounded coupling threshold is $0.01$.  The modeled matrix contains $81$ entries above that threshold, with
$$
\begin{aligned}
\max_{i,j}|S_{ij}| &= 3801961.031607896,\\
\max_i |S_{ii}| &= 3801961.031607896,\\
\tau &= 0.01.
\end{aligned}
$$
The registered break condition for this run is met only if the finite modeled-$tAI$ matrix has no $|S_{ij}|$ above the bounded-contact threshold, or if a powered survival claim is demanded without at least $40$ matched species.  The first alternative is false because $81$ entries exceed $0.01$.  The second limitation is active for any population-powered interpretation because the organism count is only $3$.  Consequently the correct readback is: nonzero bounded modeled coupling exists on the finite adaptation-weight matrix, while powered translation survival remains unavailable from this panel.

\paragraph{Layer discipline.}
The matrix is separated from direct protein claims.  A residual-to-protein association is not admissible here unless the same residual first passes through a translation readout.  In the larger matched protein-abundance mediation run, the same nine-coordinate set is used across $12$ organism panels, and each panel has rank-sufficient residualized designs with $9$ $Q$ coordinates and $24$ control columns before rank reduction.  The joined protein rows include $4042$ proteins in \emph{Bacillus subtilis}, $11661$ in \emph{Caenorhabditis elegans}, $14410$ in zebrafish, $13295$ in fly, $3739$ in \emph{E. coli} K-12 MG1655, $12418$ in chicken, $19053$ in human, $19395$ in mouse, $4129$ in \emph{Mycobacterium smegmatis}, $13322$ in rat, $5815$ in yeast, and $1656$ in \emph{Sulfolobus solfataricus}.  Those downstream rows show that a controlled $Q\to T\to P$ decomposition can be formed when protein abundance is joined, but they do not turn this three-organism modeled-$tAI$ matrix into measured ribosome occupancy.  In yeast, for example, the protein-abundance mediation fraction in the joined decomposition is $0.9491551041216973$ with total $R^{2}_{QP}=0.29837624994140827$ and direct conditioned $R^{2}_{QP\mid T}=0.0151709093608293$; that number is a statistical decomposition under modeled $tAI$, not a causal mechanism.

\paragraph{Controls and nearby negative boundaries.}
The admissible control language for this packet keeps amino-acid composition, organism, GC3, CDS length, and the $M$-density baseline as controls rather than promoted biological coordinates.  In the yeast structure-mediated abundance audit, the same nine-coordinate residual design is residualized against an intercept, $\log(\mathrm{CDS\ length})$, $20$ amino-acid composition proportions, GC3, and an $M$-density baseline.  That five-way join has $4759$ genes, with a reported CDS join hit rate $0.9156038419146592$ and an overall join hit rate $0.8818572656921754$.  Its $5'$ CDS MFE component explains only a small structure association with $S_Q$ structure $R^2=0.003774248455422586$ and permutation right-tail $p=0.027972027972027972$, while the incremental structure fraction beyond measured translation efficiency and mRNA stability is $0.0005507778820193332$ with permutation right-tail $p=0.1258741258741259$.  These scalars reinforce the local rule: bounded codon-coordinate signal may be real at a controlled readout, but structure, stability, and protein-state interpretations require their own matched contacts.

\paragraph{What the result supports.}
This finite record supports a modeled translation-readout contact for the $B^{*}_{Q6}$ residual basis.  The statement is not that a universal biological law has been found, nor that a ribosome has been observed to execute the residual coordinate.  The statement is narrower: on a three-organism modeled-$tAI$ surface with complete sense-codon weights, the nine residual coordinates produce nonzero bounded coupling entries above $0.01$.  This is sufficient to keep the $B^{*}_{Q6}$ basis alive at the modeled translation-readout boundary and to justify seeking larger, measured, organism-matched panels.

\paragraph{What the result refuses.}
No translation mechanism, folding consequence, physical admissibility, molecular function, phenotype, selection law, or universality follows from this matrix alone.  Modeled $tAI$ or $stAI$ weights are not ribosome profiling.  Matched mRNA abundance is a transcription-layer control, not a translation readout.  Protein abundance, when used, is a downstream steady-state measurement influenced by expression, turnover, localization, post-translational state, and assay bias.  The evolutionary selection-imprint audit illustrates the same refusal: with $80$ orthogroups, $240$ rows, and $3$ species, the positive selection-imprint gate gives $\Delta R^2=0.03250515015613468$, leave-clade-out $\Delta R^2=-0.0021477729853302424$, and null $95\%$ maximum $\Delta R^2=0.05427157066522958$, so the strong selection-imprint claim is not supported.  The finite coded-table phase remains separately confirmatory for lower-layer genetic-code geometry, with $40$ successes out of $40$ in the independent replication statistic and exact one-sided binomial upper-tail probability $9.094947017729282\times 10^{-13}$, but that lower-layer module success does not promote this modeled-$tAI$ matrix to protein realization.

\paragraph{Verdict.}
The chapter-level verdict is bounded positive contact at the modeled translation-readout layer and refusal of powered survival.  The finite facts are internally coherent: $61$ sense codons are weighted in each of $3$ organisms, the nine residual rows are computed, $81$ modeled entries exceed the $0.01$ coupling threshold, and the maximum absolute entry is $3801961.031607896$.  The same facts also explain the limitation: $3$ organisms are far below the stated $40$-species requirement for a powered survival claim.  The correct name for the result is therefore modeled-$tAI$ bounded coupling underpowered for measured translation survival.